\scanpage{187}\mspage{43}
% archive sheet number: 187

des $W_\infty$-fibrations faibles et a fortiori des $W_0$-fibrations faibles), mais pas \unsure{de|des} $W_0$-fibrations ; car dans le cas où $I',I$ sont $0$-connexes, il faudrait que les fibres soient $0$-connexes \inserted{(cf b) ci-dessus)}, i.e. que $u$ soit un iso !

Une propriété \unsure{sympa.} des morphismes $W_0$-lisses (resp. $W_0$-propres), c'est qu'\deleted{\unsure{elles}} ce sont \inserted{des applications} \deleted{sont} \uline{ouvertes} \inserted{(resp. fermées)} pour les espaces topologiques sous-jacents), ce qui signifie en effet que les fibres locales i.e. des flèches localisées au but \inserted{$I'_{/\unsure{i'|u i'}}\longrightarrow I_{/\unsure{i|i'}}$} (resp colocales i.e. les flèches colocalisées au but $\,{}_{i'\backslash}I'\longrightarrow{}_{i\backslash}I$) sont surjectives, i.e. à fibres non-vides (ce qui est impliqué par la condition que celles-ci soient $0$-connexes).

Dans tout ça je n'ai \deleted{pas} encore guère parlé des flèches $W_0$-asphériques resp. co-asphériques, qui pourtant sont ici ma motivation \unsure{principale} pour introduire $W_0$ !

\uline{Proposition 3.2.} Conditions équivalentes sur la flèche $u:I'\longrightarrow I$ de $\mathrm{Cat}$ :

a) Pour tte $\mathcal{U}$-catégorie \inserted{$M$}, et tout foncteur $F:I\longrightarrow M$, la flèche canonique (sous réserve que les \deleted{\unsure{limites} \unsure{ind}} \inserted{petites} limites inductives sont \unsure{repr.|représ.} dans $M$)

\[
\underset{I'}{\underset{\longrightarrow}{\lim}}\,u^*\unsure{(}\deleted{\unsure{F}}F\unsure{)|,}\ \longrightarrow\ \underset{I}{\underset{\longrightarrow}{\lim}}\,F \tag*{\unsure{(3.2.1.)|(3.2,1.)}}
\]
